\section{Application summary (550 words)}

Understanding how groups coordinate and compete --- from autonomous vehicle fleets to crowd movements --- is fundamental to many real-world challenges. Current mathematical methods for analysing these collective systems have a critical limitation: they examine behaviour at a single spatial scale, missing the hierarchical organisation (individual behaviours $\rightarrow$ group tactics $\rightarrow$ overall strategy) that defines most practical applications. This research scales a validated multi-scale topological data analysis framework for competitive spatial systems from proof-of-concept to full-season analysis, using professional football as a rigorous testbed where high-quality spatial data enables thorough method development.

\paragraph{The challenge.}
Topological Data Analysis (TDA) has proven effective for understanding the shape of complex data, but existing applications examine systems at single scales, in non-competitive scenarios, or with static data. A standard Vietoris--Rips filtration applied to a multi-agent point cloud conflates topological features from different organisational levels --- individual player interactions, tactical group structures, and team-wide organisation --- into a single persistence diagram. Separating these levels is a prerequisite for meaningful analysis of hierarchical competitive systems, yet no established method exists for doing so in a principled, domain-informed way.

\paragraph{Established framework.}
We have developed and validated a multi-scale persistent homology framework (Brown et al., 2026, ArXiv preprint) addressing this through two methodological contributions: domain-informed hierarchical clustering to decompose point clouds by organisational level, and adaptive filtration ensuring consistent $H_1$ (loop) detection across all scales. Validation across 11 professional football matches (1{,}650 analysis frames) confirms three robust $H_0$ analysis regimes (individual, tactical, team) with stability scores exceeding 0.88, and two $H_1$ regimes capturing complementary structural information --- individual-scale topology reveals frequent, transient loops (97\% presence rate) whilst tactical-scale topology reveals rarer but geometrically distinct loops (20\% presence rate). Scale complementarity is quantified (Spearman $\rho = 0.25$), confirming that neither scale subsumes the other. Real event correlation across 10 matches (104{,}722 event--topology pairs) demonstrates that topological features respond coherently to match dynamics. Temporal persistence dynamics are match-specific rather than universal --- unsurprising to practitioners --- but the framework provides the first tools to quantify this variation systematically.

\paragraph{Project objectives.}
Building on this validated foundation, three research objectives are pursued over 12 months, with a fourth strategic output: (1)~Scale the analysis pipeline to a full Championship season ($\sim$540 matches), stress-testing the framework and establishing population-level topological statistics; (2)~Characterise pattern formation by developing baseline topological fingerprints for different tactical systems, quantifying formation times and distinguishing between tactical configurations in persistence space; (3)~Develop persistence landscape methods for temporal dynamics, classifying stability regimes and characterising transitions within and across matches; (4)~Synthesise full-season results into a Standard Grant application extending the framework to broader competitive systems.

\paragraph{Broader impact.}
While validated using football data --- ideal for high-frequency spatial measurements on a well-structured competitive system --- the framework applies to any bounded competitive multi-agent system with tracking data. Applications include autonomous systems coordination, competitive crowd dynamics at transport hubs, and ecological predator--prey monitoring. The mathematical contribution sits within computational topology: practical methods for multi-scale persistent homology on hierarchical, competitive, time-evolving point clouds.

\paragraph{Pathway to impact.}
The ArXiv preprint establishes methodological validity. This Small Grant generates the full-season evidence and persistence landscape theory needed for a Standard Grant application within 6--8 months of completion, with clear industry partnership routes through established Championship club relationships and secured data access via the Swansea City AFC--StatsBomb agreement.
